\documentclass[10pt,a4paper]{amsart}
\usepackage[margin=19mm]{geometry}
\usepackage{amsmath,amssymb,mathtools}
\usepackage[scr=rsfs,cal=euler]{mathalfa}
\usepackage{xfrac}
\usepackage[unicode,hidelinks,bookmarks=false]{hyperref}
\newcommand{\ie}{i.e.\ }
\newcommand{\cf}{cf.\ }
\newcommand{\resp}{respectively\ }
\newcommand{\Subsection}{\subsection}
\providecommand{\nobreakdash}{\nobreak\hbox{-}}
\setlength{\emergencystretch}{2em}
\multlinegap0pt
\usepackage{bm}
\usepackage{bbm}
\usepackage{dsfont}
\usepackage{xspace}
\usepackage{cancel}
\usepackage{mathtools}
\usepackage{amsthm}
\makeatletter
\@ifundefined{definition}{\newtheorem{definition}{Definition}}{}
\@ifundefined{theorem}{\newtheorem{theorem}[definition]{Theorem}}{}
\makeatother
\title[Common Fixed Point Theorems Of Weakly Compatible Maps Satisf]{Common Fixed Point Theorems Of Weakly Compatible Maps Satisfying (f,g)-Weakly Contractive Condition And Invariant Approximation Results}
\author{Babu G.V.R., Ratna Babu D, Alemayehu Negash}
\date{Source: arXiv:2510.16268v1 (CC BY 4.0)}
\begin{document}
\maketitle

\noindent\emph{Source and licence.} Common Fixed Point Theorems Of Weakly Compatible Maps Satisfying (f,g)-Weakly Contractive Condition And Invariant Approximation Results. Version arXiv:2510.16268v1 (CC BY 4.0), distributed under the Creative Commons Attribution 4.0 International licence, as recorded in the arXiv licence metadata for this exact version and shown on the abstract page https://arxiv.org/abs/2510.16268v1. Full rights notice: licenses/arxiv-2510.16268-CC-BY-4.0.txt.\par\smallskip

\noindent\textit{Editorial scope.} Counted unit: the Theorem whose source label is \texttt{None} in the article (source0.tex lines 274--276, source label \texttt{None}), delivered together with the article's own complete original proof of it (source0.tex lines 278--509, one \texttt{proof} environment of 232 lines). No mathematics is written, repaired or re-derived: statement and proof are the article's own text line for line. Required context retained uncounted: none. Every definition, display and label of those lines is retained unchanged. Omitted from this package and not redistributed: 1--273, 277--277, 510--886. Nothing inside a retained excerpt is omitted or altered: every retained body is delivered exactly as the article's lines are written, so no omission receipt of any kind accompanies this package. Citations: the retained excerpts carry no citation command at all, so no bibliography entry of the article is transcribed and no reference is redistributed. Reference adaptation: none was needed and none was made; every reference inside the delivered unit resolves inside it, so no unresolved reference or citation is produced. Printed slips kept literal: no printed word, symbol, hypothesis or number of the counted statement or of its proof was altered. Layout and class: the article's own class is \texttt{article} with packages including amsmath, amssymb, graphicx, amsthm, mathrsfs, mathptmx, mathastext, mathtools, inputenc, enumitem, hyperref, geometry; the wrapper is the lane's pinned \texttt{amsart} editorial wrapper at 10pt on A4, which restates the theorem-like environments the retained text needs and adds the article's own macro definitions that the retained text needs (none), with extra wrapper packages (bm, bbm, dsfont, xspace, cancel, mathtools). The delivered numbering is the unit's own. Assets: no retained span contains a figure environment, a graphics-inclusion command, a PDF-page inclusion, a table or a TikZ picture; no image file is copied into this package, the package declares no asset dependency and no image file is redistributed with it. Licence: version arXiv:2510.16268v1 of this article is distributed under the Creative Commons Attribution 4.0 International licence, as recorded in the arXiv licence metadata for this exact version and shown on the abstract page https://arxiv.org/abs/2510.16268v1. A full rights notice is delivered at licenses/arxiv-2510.16268-CC-BY-4.0.txt. Characters: every retained excerpt is pure ASCII, so the delivered engine, XeLaTeX with its default Latin Modern text font, reports no missing character.
\begin{theorem}
Let $(X,d)$ be a metric space and let $T$ be a $(f,g)-$ weakly contractive mapping. If the pairs $(T,f)$ and $(T,g)$ are weakly compatible and $\overline{T(X)}\subseteq f(X)$, $\overline{T(X)}\subseteq g(X)$ and $\overline{T(X)}$ is a complete subspace of $X$, then $T$, $f$ and $g$ have a unique common fixed point in $X$.
\end{theorem}

\begin{proof}
Let $x_{0}\in X.$ Since $\overline{T(X)}\subseteq f(X)$, $T(X)\subseteq f(X).$ So there exists $x_{1}\in X$ such that $Tx_{0}=fx_{1}.$
Since $\overline{T(X)}\subseteq g(X)$, $T(X)\subseteq g(X).$ So there exists $x_{2}\in X$ such that $Tx_{1}=gx_{2}.$
On continuing this process, inductively we get a sequence $\{x_{n}\}$ in $X$ such that $y_{2n}=T(x_{2n})=f(x_{2n+1})$ and $y_{2n+1}=T(x_{2n+1})=g(x_{2n+2})$, $n=0,1,2,\ldots.$
We now show that the sequence $\{d(y_{n},y_{n+1})\}$ is a decreasing sequence. Now,
\begin{align*}
d(y_{2n},y_{2n+1}) &= d(T(x_{2n}),T(x_{2n+1})) \\
&\leq \min\{d(f(x_{2n}),g(x_{2n+1}))-\psi(d(f(x_{2n}),g(x_{2n+1}))), \\
&\quad d(g(x_{2n}),f(x_{2n+1}))-\psi(d(g(x_{2n}),f(x_{2n+1})))) \\
&\leq d(g(x_{2n}),f(x_{2n+1}))-\psi(d(g(x_{2n}),f(x_{2n+1}))) \\
&= d(y_{2n-1},y_{2n})-\psi(d(y_{2n-1},y_{2n})) \\
&< d(y_{2n-1},y_{2n}).
\end{align*}
Therefore
\begin{equation}
d(y_{2n},y_{2n+1})<d(y_{2n-1},y_{2n}).
\end{equation}
%(2.1.1)
Similarly, it is easy to see that
\begin{equation}
d(y_{2n+1},y_{2n+2})<d(y_{2n},y_{2n+1}).
\end{equation}
%(2.1.2)
Thus from (2.1.1) and (2.1.2) it follows that

\[d(y_{n},y_{n+1})<d(y_{n-1},y_{n}),\ n=0,1,2,\ldots.\]
Therefore $\{d(y_{n},y_{n+1})\}$ is a strictly decreasing sequence of nonnegative reals and hence convergent and let its limit be $l\geq 0$.
If $l>0$ then $l=\inf\limits_{n}d(y_{n},y_{n+1})$.
Thus we have $l\leq d(y_{n},y_{n+1})$ for all $n$.\\
Now,
\begin{align}
d(y_{n},y_{n+1}) &= d(T(x_{n}),T(x_{n+1})) \nonumber \\
&\leq \min\{d(f(x_{n}),g(x_{n+1}))-\psi(d(f(x_{n}),g(x_{n+1}))), \nonumber \\
&\quad d(g(x_{n}),f(x_{n+1}))-\psi(d(g(x_{n}),f(x_{n+1})))\}.
\end{align}
%(2.1.3)
If $n$ is odd then from (2.1.3), we get
\begin{align}
d(y_{n},y_{n+1}) &= d(T(x_{n}),T(x_{n+1})) \nonumber \\
&\leq d(f(x_{n}),g(x_{n+1}))-\psi(d(f(x_{n}),g(x_{n+1}))) \nonumber \\
&= d(y_{n-1},y_{n})-\psi(d(y_{n-1},y_{n})).\nonumber
\end{align}
%(2.1.4)
Therefore
\begin{equation}
    d(y_{n},y_{n+1})\leq d(y_{n-1},y_{n})-\psi(d(y_{n-1},y_{n})).
\end{equation}
If $n$ is even then from (2.1.3), we get
\begin{align}
d(y_{n},y_{n+1}) &= d(T(x_{n}),T(x_{n+1})) \nonumber \\
&\leq d(g(x_{n}),f(x_{n+1}))-\psi(d(g(x_{n}),f(x_{n+1}))) \nonumber \\
&= d(y_{n-1},y_{n})-\psi(d(y_{n-1},y_{n})).\nonumber
\end{align}
%(2.1.5)
Therefore
\begin{equation}
    d(y_{n},y_{n+1})\leq d(y_{n-1},y_{n})-\psi(d(y_{n-1},y_{n})).
\end{equation}
We have, from (2.1.4) and (2.1.5)
\begin{equation}
d(y_{n},y_{n+1})\leq d(y_{n-1},y_{n})-\psi(d(y_{n-1},y_{n}))\ \ \text{for all}\ \ n.
\end{equation}
%(2.1.6)
On taking limit as $n\to\infty$ in (2.1.6), we get
\[
\lim\limits_{n\to\infty}d(y_{n},y_{n+1})\leq\lim\limits_{n\to\infty}d(y_{n-1},y_{n})-\psi(\lim\limits_{n\to\infty}d(y_{n-1},y_{n})).
\]
Therefore $l\leq l-\psi(l)<l$, a contradiction.
Hence, $l=0$.\\
Thus it follows that
\[
\lim_{n\to\infty}d(y_n, y_{n+1}) = 0.
\]
\textbf{Case(i):}  $y_n = y_{n+1}$ for some $n$.\\
If $n$ is even, we write $n = 2m$ (say), $m \in Z_+$.
Thus, from (1.8.1) it follows that
\[
y_{2m} = y_{2m+k} \quad \text{for all } k = 0, 1, 2, \ldots.
\]
If $n$ is odd then $n = 2m + 1$ (say), $m \in Z_+$.
From (1.8.1), we get
\[
y_{2m+1} = y_{2m+1+k} \quad \text{for all } k = 0, 1, 2, \ldots.
\]
Therefore
\[
y_n = y_{n+k} \quad \text{for all } k = 0, 1, 2, \ldots.
\]
Thus it follows that $\{y_m\}_{m\geq n}$ is a constant sequence and hence it is a Cauchy sequence.\\
\textbf{Case(ii):}  $y_n \neq y_{n+1}$ for all $n$.\\
We now show that $\{y_n\}$ is a Cauchy sequence in $T(X)$.
It is sufficient to show that the sequence $\{y_{2n}\}$ is a Cauchy sequence.
Suppose $\{y_n\}$ is not Cauchy. Then there exists $\epsilon > 0$ for all $k \in \mathbb{N}$ there exist $\{m_k\}$, $\{n_k\}$ with $n_k > m_k > k$ such that
\begin{equation}
d(y_{2m_k}, y_{2n_k}) \geq \epsilon \quad \text{and} \quad d(y_{2m_k}, y_{2n_k-2}) < \epsilon.
\end{equation}
%(2.1.7)
Now,
\begin{equation}
d(y_{2n_k}, y_{2m_k}) \leq d(y_{2n_k}, y_{2n_k-1}) + d(y_{2n_k-1}, y_{2n_k-2}) + d(y_{2n_k-2}, y_{2m_k}).
\end{equation}
%(2.1.8)
On taking limits as $k \to \infty$ in (2.1.8) and using (2.1.7), we get
\[
\epsilon \leq \lim_{k\to\infty}d(y_{2n_k}, y_{2m_k}) + \lim_{k\to\infty}d(y_{2n_k-1}, y_{2n_k-2}) + \lim_{k\to\infty}d(y_{2n_k-2}, y_{2m_k}) \leq \epsilon.
\]
Therefore
\begin{equation}
\lim_{k\to\infty}d(y_{2n_k}, y_{2m_k}) = \epsilon.
\end{equation}
%(2.1.9)
We now show that $\lim\limits_{k\to\infty}d(y_{2m_k-1}, y_{2n_k}) = \epsilon$.
Now,
\begin{equation}
d(y_{2m_k-1}, y_{2n_k}) \leq d(y_{2m_k-1}, y_{2n_k-1}) + d(y_{2n_k-1}, y_{2n_k}).
\end{equation}
%(2.1.10)
On taking limits as $k\to\infty$ in (2.1.10), we get
\begin{align*}
\lim_{k\to\infty}d(y_{2m_{k}-1},y_{2n_{k}}) &\leq \lim_{k\to\infty}d(y_{2m_{k}-1},y_{2n_{k}-1})+\lim_{k\to\infty}d(y_{2n_{k}-1},y_{2n_{k}}) \\
&\leq \epsilon.
\end{align*}
Therefore
\begin{equation}
\lim_{k\to\infty}d(y_{2m_{k}-1},y_{2n_{k}}) \leq \epsilon.
\end{equation}
%(2.1.11)
Now,
\begin{align}
\epsilon &\leq d(y_{2m_{k}-1},y_{2n_{k}-1}) \nonumber\\
&\leq d(y_{2m_{k}-1},y_{2n_{k}})+d(y_{2n_{k}},y_{2n_{k}-1})
\end{align}
%(2.1.12)
On taking limits as $k\to\infty$ in (2.1.12), we get
\begin{equation}
\epsilon\leq\lim_{k\to\infty}d(y_{2m_{k}-1},y_{2n_{k}}).
\end{equation}
%(2.1.13)
We have, from (2.1.11) and (2.1.13)
\begin{equation}
\lim_{k\to\infty}d(y_{2m_{k}-1},y_{2n_{k}})=\epsilon.
\end{equation}
%(2.1.14)
Thus by using (2.1.9) and (2.1.14), we get
\begin{align*}
\epsilon &= d(y_{2m_{k}},y_{2n_{k}}) \\
&\leq d(y_{2m_{k}},y_{2n_{k}+1})+d(y_{2n_{k}+1},y_{2n_{k}}) \\
&= d(T(x_{2m_{k}}),T(x_{2n_{k}+1}))+d(y_{2n_{k}+1},y_{2n_{k}}) \\
&\leq \min\{d(f(x_{2m_{k}}),g(x_{2n_{k}+1}))-\psi(d(f(x_{2m_{k}}),g(x_{2n_{k}+1}))), \\
&\quad d(g(x_{2m_{k}}),f(x_{2n_{k}+1}))-\psi(d(g(x_{2m_{k}}),f(x_{2n_{k}+1})))\} \\
&\quad +d(y_{2n_{k}+1},y_{2n_{k}}) \\
&\leq d(g(x_{2m_{k}}),f(x_{2n_{k}+1}))-\psi(d(g(x_{2m_{k}}),f(x_{2n_{k}+1}))) \\
&\quad +d(y_{2n_{k}+1},y_{2n_{k}}) \\
&= d(y_{2m_{k}-1},y_{2n_{k}})-\psi(d(y_{2m_{k}-1},y_{2n_{k}})))+d(y_{2n_{k}+1},y_{2n_{k}}).
\end{align*}
Therefore
\begin{equation}
\epsilon \leq d(y_{2m_{k}-1},y_{2n_{k}})-\psi(d(y_{2m_{k}-1},y_{2n_{k}}))+d(y_{2n_{k}+1},y_{2n_{k}}).
\end{equation}
%(2.1.15)
On taking limits as $k\to\infty$ in (2.1.15), we get
\begin{align*}
\epsilon &\leq \lim_{k\to\infty}d(y_{2m_{k}-1},y_{2n_{k}})-\lim_{k\to\infty}\psi(d(y_{2m_{k}-1},y_{2n_{k}}))+\lim_{k\to\infty}d(y_{2n_{k}+1},y_{2n_{k}}) \\
&= \epsilon-\psi(\epsilon)<\epsilon,
\end{align*}
a contradiction.
Therefore $\{y_{2n}\}$ is a Cauchy sequence in $T(X)$ and hence $\{y_{n}\}$ is Cauchy in $T(X)$.\\
Since $T(X)\subseteq\overline{T(X)}$, we have $\{y_{n}\}$ is Cauchy in $\overline{T(X)}$ and $\overline{T(X)}$ is complete there exists $z\in T(X)$ such that $y_{n}\to z$ as $n\to\infty$.\\
Now,
\[
\lim_{n\to\infty}y_{2n}=\lim_{n\to\infty}T(x_{2n})=\lim_{n\to\infty}f(x_{2n+1})=z
\]
and
\[
\lim_{n\to\infty}y_{2n+1}=\lim_{n\to\infty}T(x_{2n+1})=\lim_{n\to\infty}g(x_{2n+2})=z.
\]
Since $\overline{T(X)}\subseteq f(X)$ and $z\in\overline{T(X)}$, there exists $u\in X$ such that $z=f(u)$ and since $T(X)\subseteq g(X)$ and $z\in T(X)$, there exists $v\in X$ such that $z=g(v)$.\\
We show now that $z=T(u)$ and $z=T(v)$.\\
Now,
\begin{align}
d(Tu,T(x_{2n+2})) &\leq \min\{d(fu,g(x_{2n+2}))-\psi(d(fu,g(x_{2n+2}))), \nonumber \\
&\quad d(gu,f(x_{2n+2}))-\psi(d(gu,f(x_{2n+2})))\} \nonumber \\
&\leq d(fu,g(x_{2n+2}))-\psi(d(fu,g(x_{2n+2}))).
\end{align}
%(2.1.16)
On taking limits as $n\to\infty$ in (2.1.16), we get
\[
d(Tu,z) \leq d(z,z)-\psi(d(z,z)) = 0.
\]
Thus it follows that $Tu=z$.\\
Now,
\begin{align}
d(Tv,T(x_{2n+1})) &\leq \min\{d(fv,g(x_{2n+1}))-\psi(d(fv,g(x_{2n+1}))), \nonumber \\
&\quad d(gv,f(x_{2n+1}))-\psi(d(gv,f(x_{2n+1})))\} \nonumber \\
&\leq d(gv,f(x_{2n+1}))-\psi(d(gv,f(x_{2n+1}))).
\end{align}
%(2.1.17)
On taking limits as $n\to\infty$ in (2.1.17), we get
\[
d(Tv,z) \leq d(z,z)-\psi(d(z,z)) = 0.
\]
Thus it follows that $Tv=z$.\\
Since $(T,f)$ and $(T,g)$ are weakly compatible and since $Tu=fu$ and $Tv=gv$ we have
\[
fz=f(Tu)=T(fu)=Tz=T(gv)=g(Tv)=gz.
\]
This shows that $z$ is a coincidence point of $T,f$ and $g$.\\
We now show that $Tz=z$.\\
Consider
\begin{align*}
d(Tz,z) &= d(Tz,Tu) \\
&\leq \min\{d(fz,gu)-\psi(d(fz,gu)),d(gz,fu)-\psi(d(gz,fu))\} \\
&\leq d(gz,fu)-\psi(d(gz,fu)) \\
&= d(Tz,z)-\psi(d(Tz,z))
\end{align*}
Therefore $d(Tz,z)\leq d(Tz,z)-\psi(d(Tz,z))$, which implies that $\psi(d(Tz,z))\leq 0$. This shows that $d(Tz,z)=0$ and hence $Tz=z$.\\
Therefore $z$ is a common fixed point of $T,f$ and $g$.

\subsection*{Uniqueness of $z$:}

Let $t$ be another common fixed point of $T,f$ and $g$ with $t\neq z$.
Then
\begin{align*}
d(t,z) &= d(Tt,Tz) \\
&\leq \min\{d(ft,gz)-\psi(d(ft,gz)),d(gt,fz)-\psi(d(gt,fz))\} \\
&\leq d(ft,gz)-\psi(d(ft,gz)) \\
&= d(t,z)-\psi(d(t,z)).
\end{align*}
Thus it follows that $\psi(d(t,z))\leq 0$ so that $d(t,z)=0$ and hence $t=z$.\\
Therefore $z$ is the unique common fixed point of $T,f$ and $g$.\\
Hence the theorem follows.
\end{proof}

\end{document}
